\section*{الفصل \ref{ch:b2:asexual-reproduction} --- التكاثر اللاجنسي والاستنساخ عند النباتات}

\begin{solution}{exo:b2:asexual-reproduction:1}
\omterm{def:b2:asexual-reproduction:clone}{النسيلة}: الخلف المتطابق وراثيًا لأب واحد بالتقسم
الخيطي. \omterm{def:b2:asexual-reproduction:clone}{والجينيت}: الفرد الوراثي، أي كل ما نما من لاقحة
واحدة. \omterm{def:b2:asexual-reproduction:clone}{والراميت}: وحدة مستقلة فيزيولوجيًا من \omterm{def:b2:asexual-reproduction:clone}{جينيت}. والأعضاء:
المدادات (الفراولة)، والجذامير (النجيل)، والدرنات
(البطاطا)، والبصلات (البصل)، والخلفات (الحور الرجراج)، ونبيتات الورقة
(\emph{Kalanchoe}).
\end{solution}

\begin{solution}{exo:b2:asexual-reproduction:2}
يجب أن يتحاذى كامبيوما الطُّعم والأصل الوعائيان ويتلامسا؛ ويجسر
الدشبذ القطع ويمايز عبره خشبًا ولحاء. ويوفّر الأصل الجذور (الماء،
والمعادن، والقوة، واحتمال التربة، وضبط القياس)، ويوفّر الطُّعم
الفرع المثمر (أي الصنف). ويبقي كل جزء مورثومه، ولا تندمج الخلايا
أبدًا: أي نمطان وراثيان جنبًا إلى جنب، وهذه شيمرا؛ أما الهجين فكان
ليحمل خليطًا من المورثومين في كل خلية.
\end{solution}

\begin{solution}{exo:b2:asexual-reproduction:3}
\omterm{def:b2:asexual-reproduction:clone}{التكاثر الخضري}: جزء من الجسم يتجذر نباتًا جديدًا.
\omterm{def:b2:asexual-reproduction:apomixis}{والتوالد العذري البذري}: \omterm{def:b2:angiosperm-reproduction:seed}{بذرة} تتكون بلا إلقاح، وجنينها \omterm{def:b2:asexual-reproduction:clone}{نسيلة} من
الأم. \omterm{def:b2:asexual-reproduction:apomixis}{والتوالد العذري}: حيوان ينمو من بويضة غير ملقحة.
ولا يبقي \omterm{def:b2:angiosperm-reproduction:seed}{البذرة} --- بقشرتها وسكونها وانتشارها --- إلا \omterm{def:b2:asexual-reproduction:apomixis}{التوالد
العذري البذري}.
\end{solution}

\begin{solution}{exo:b2:asexual-reproduction:4}
عقّم المستنبت؛ وضعه على آغار فيه سكر ومعادن
وفيتامينات وأوكسين وسيتوكينين متوازنان (فيتكون الدشبذ)؛ ثم ارفع
نسبة السيتوكينين إلى الأوكسين لاستحثاث الفروع؛ وأعد زراعة الفروع
كل بضعة أسابيع؛ وانقلها إلى وسط غني بالأوكسين للتجذير؛ ثم افطم
النبيتات في رطوبة عالية ثم في التربة.
\end{solution}

\begin{solution}{exo:b2:asexual-reproduction:5}
يُستبدل بكل نبات $1 + 24 = 25$ نباتًا: فبعد 4 سنوات يكون $25^{4} =
\num{390625}$ نباتًا، تحتاج إلى \qty{39000}{m^2}، أي نحو أربعة هكتارات.
والمكان والضوء والمغذيات، وقصر مدى المداد، وموت النباتات تحت
الازدحام، توقف الطور الأسّي في غضون سنتين؛ ثم تتقدم \omterm{def:b2:asexual-reproduction:clone}{النسيلة} جبهةً.
\end{solution}

\begin{solution}{exo:b2:asexual-reproduction:6}
$\num{14000}/130 \approx 110$ جيلًا من الجذوع. ونصف قطر قرص
مساحته \qty{4.3e5}{m^2}: $\sqrt{4.3\times 10^{5}/\pi} = \qty{370}{m}$؛
فالتقدم $370/\num{14000} = \qty{0.026}{m}$ في السنة، أي دون ثلاثة
سنتيمترات.
\end{solution}

\begin{solution}{exo:b2:asexual-reproduction:7}
$p_t = 2^{t}p_0/(1 + (2^{t} - 1)p_0)$: فعند $t = 5$: $0.0032$؛ وعند $t = 10$:
$0.1024/1.1023 = 0.093$؛ وعند $t = 15$: $3.28/4.28 = 0.77$. ولا تنتشر
أبدًا إذا كانت لياقتها دون نصف لياقة السلالة الجنسية: فلياقة
نسبية $w$ تعطي $q' = 2wq$، وهذا يتناقص عند $w <
\tfrac12$.
\end{solution}

\begin{solution}{exo:b2:asexual-reproduction:8}
$10 : 1$: جذور؛ و$1 : 1$: دشبذ غير متمايز؛ و$1 : 10$: فروع.
والعقلة لها فرع سلفًا وتحتاج إلى جذور، والأوكسين يستحثها
عند القاعدة المقطوعة.
\end{solution}

\begin{solution}{exo:b2:asexual-reproduction:9}
1 ثم 10 ثم 100 ثم \num{1000} ثم \num{10000} كغ: فبعد الموسم الثالث
\qty{1}{t} فقط، وبعد الرابع \qty{10}{t}: أي أربعة مواسم. \omterm{def:b2:angiosperm-reproduction:seed}{والبذرة}
الحقيقية تعطي آلاف البذور في النبات لكن البطاطا رباعية صيغة متغايرة
اللواقح، فتنفصل البادرات ولا تكون أي منها الصنف.
\end{solution}

\begin{solution}{exo:b2:asexual-reproduction:10}
$e^{-5} = 0.0067$: أي 670 فردًا خاليًا من الطفرات في الجمهرة
الكبيرة، ونحو 7 في الصغيرة. والسبعة قد يخفقون كلهم في
ترك خلف بالمصادفة خلال أجيال قليلة (فاحتمال ألا يتكاثر أي
منهم من رتبة $e^{-7}$ في الجيل)، وعندها يضيع أفضل صنف إلى الأبد
ويصير الذي يليه هو الأفضل: أي نقرة من السقّاطة، تتكرر حتى
يصير العبء لا يُحتمل. أما مع 670 فالضياع مستحيل عمليًا
--- إلى أن يقلّص عنق زجاجة العدد $N$ أو يرتفع \omterm{thm:b2:genome-diversification:rate}{معدل الطفر}؛
فالسقّاطة لا تدور إلا في اتجاه واحد.
\end{solution}

\begin{solution}{exo:b2:asexual-reproduction:11}
في \omterm{def:b2:asexual-reproduction:clone}{النسيلة} يكون كل مضيف قابلًا للإصابة: فتنتشر السلالة
بلا كابح بسرعة سفر الأبواغ. أما في الجمهرة الجنسية
\omterm{def:b2:meiosis-heredity:vocabulary}{فالنمط الوراثي} الذي تتغلب عليه واحد من $2^{10} = 1024$ تركيبًا،
موجود في نحو جزء من ألف من النباتات، مبعثر بين جيران
مقاومين، فتجوع الوبيئة. وكان اللَّمبر \omterm{def:b2:asexual-reproduction:clone}{نسيلة} واحدة
مزروعة في معظم بلد؛ فوجدت اللفحة كل نبات
مفتوحًا، ولم يكن ثمة تنوع يُربّى منه حتى استُوردت أصناف أخرى.
\end{solution}

\begin{solution}{exo:b2:asexual-reproduction:12}
\omterm{prop:b2:asexual-reproduction:whysex}{سقّاطة مولر}، واجتماع الطفرات المفضّلة، والملكة الحمراء: كل
منها يعطي الجنس فائدة قد تفوق عامل اثنين في عالم كبير
مليء بالطفيليات متغير. والدولابيات البدلوئيدية هي الحالة
المحرجة: أربعون مليون سنة بلا ذكور وبلا انهيار سقّاطة.
ويبدو أنها تفلت بمقاومة شديدة للجفاف
(وهي تقتل طفيلياتها)، وبالتقاط DNA غريب، وبإصلاح
DNA فعال إلى حد غير معتاد --- وهي استثناءات تبين ما تقتضيه
القاعدة عادةً.
\end{solution}

\begin{solution}{pb:b2:asexual-reproduction:1}
\textbf{1.} $4\times 10^{4}$ نبات.
\textbf{2.} $1 + 6\times 2 = 13$.
\textbf{3.} \num{1300}؛ ثم \num{16900}؛ ثم \num{219700}.
\textbf{4.} $13^{t} = 400$: ومنه $t = \ln 400/\ln 13 = 5.99/2.56 = 2.3$
سنة --- أي يمتلئ خلال السنة الثالثة.
\textbf{5.} نصف قطر قرص مساحته هكتار هو $\sqrt{10^{4}/\pi} = \qty{56}{m}$؛
وعند \qty{0.5}{m} في السنة يلزم \num{113} سنة.
\textbf{6.} مئة مؤسِّس مبعثر تملأ الحقل أسّيًا
من مئة جبهة في سنتين؛ أما المؤسِّس الواحد فيملأ جواره
في سنتين ثم يزحف قرنًا.
\textbf{7.} $5^{n} \ge 10^{6}$: ومنه $n \ge 8.6$، أي تسع إعادات زراعة
($1.95\times 10^{6}$ فرع)، و36~أسبوعًا.
\textbf{8.} $0.8\times 1.95\times 10^{6} = 1.6\times 10^{6}$ نبات.
\textbf{9.} $0.98$؛ و$0.98^{10} = 0.82$.
\textbf{10.} $10\times 0.98 = 9.8$ خطًا نظيفًا منتظرًا؛ ومن عشر
عقل، واحدة.
\textbf{11.} ينتقل الفيروس عبر الوصلات الهيولية واللحاء،
والقبة القمية لا وصل وعائي لها وتنقسم أسرع
مما ينتشر الفيروس، فيكون العشر الخارجي من المليمتر
غير مصاب عادةً.
\textbf{12.} المختبر: نباتات موحدة خالية من الفيروس حالًا، لكنها
مكلفة، وأي \omterm{def:b2:genome-diversification:mutation}{طفرة} أو تلوث في المزرعة يُضاعف
مليون مرة. والمدادات: مجانية وفي الموقع، لكن ثلاثة
مواسم تضيع وكل فيروس في النباتات الأم يُحمل معها.
\textbf{13.} عدد المتوالدات عذريًا $pN$ وتعطي كل منها $F$ \omterm{def:b2:angiosperm-reproduction:seed}{بذرة}؛ والجنسيات
$(1 - p)N$ وتعطي كل منها $F/2$ \omterm{def:b2:angiosperm-reproduction:seed}{بذرة} (إذ يذهب نصف جهدها لقاحًا لا
يصنع بذورًا)؛ فنصيب المتوالدة عذريًا من البذور هو $Fp/(Fp + F(1 - p)/2)
= 2p/(1 + p)$.
\textbf{14.} $p_5 = 0.032/1.031 = 0.031$؛ و$p_{10} = 1.024/2.023 =
0.51$: فتتجاوز النصف عند الجيل 10.
\textbf{15.} $p' = 2p(1 - cp)/\bigl(2p(1 - cp) + (1 - p)\bigr)$.
\textbf{16.} يكون $p' = p$ حين $2(1 - cp^*) = 1$: أي $p^* = 1/2c = 0.625$.
ودون $p^*$ يكون $2(1 - cp) > 1$ فيرتفع $p$؛ وفوقه يهبط: فهو مستقر.
\textbf{17.} يكون $p^* \ge 1$ حين $c \le \tfrac12$: فتستولي المتوالدة عذريًا
ما لم يكلفها الطفيلي أكثر من نصف إنتاجها حين تشيع.
\textbf{18.} عند $c = 0.6$: $p^* = 0.83$. ولإبقاء الجنس في الجمهرة
على الطفيلي أن يقطع، عند تواتر عالٍ للمتوالدة عذريًا، من لياقة
\omterm{def:b2:asexual-reproduction:clone}{النسيلة} أكثر من الميزة المضاعفة --- أي على الملكة الحمراء أن
تركض بسرعة.
\textbf{19.} ثلاثة طقوم صبغية لا تستطيع التزاوج في \omterm{def:b2:meiosis-heredity:meiosis}{التقسم المنصّف}،
فتكون أعراسها غير متوازنة؛ فهي محبوسة خارج إعادة التركيب،
وتعمل عليها \omterm{prop:b2:asexual-reproduction:whysex}{سقّاطة مولر} والطفيليات بلا معارض --- ولهذا
كانت سلالات الهندباء المتوالدة عذريًا حديثة وتُعاد صناعتها
من أسلاف جنسية.
\textbf{20.} $10^{-9}\times 4.8\times 10^{8} = 0.48$ في الانقسام.
\textbf{21.} $2\times \num{14000}\times 20 = 5.6\times 10^{5}$
انقسام؛ ومنه $0.48\times 5.6\times 10^{5} = 2.7\times 10^{5}$
\omterm{def:b2:genome-diversification:mutation}{طفرة}.
\textbf{22.} $2.7\times 10^{5}/4.8\times 10^{8} = 5.6\times 10^{-4}$ من
المورثوم؛ أي نحو \num{5400} فرق مرمّز.
\textbf{23.} التنافس بين سلالات الخلايا في المرستيم (فالسلالة
الطافرة الأبطأ انقسامًا تُزاح) وبين الراميتات (فالجذع الذي
يحمل \omterm{def:b2:genome-diversification:mutation}{طفرة} سيئة يخلّف خلفات أقل). وهو أضعف لأن معظم الطفرات
الجسدية متغايرة اللواقح ومقنّعة، ولا شيء يكشفها:
فلا \omterm{def:b2:meiosis-heredity:meiosis}{تقسم منصّف} يجعلها متماثلة اللواقح ولا طور لاقحة وحيدة
الخلية يختبر المورثوم كله في خلية واحدة.
\textbf{24.} ليست موحدة: فجذوعها تختلف في مئات الآلاف
من المواقع. لكن الفروق محايدة مقنّعة في معظمها، \omterm{def:b2:asexual-reproduction:clone}{فالنسيلة}
\omterm{def:b2:meiosis-heredity:vocabulary}{نمط وراثي} واحد فعليًا --- أي فسيفساء من متغيرات جسدية
\omterm{def:b2:asexual-reproduction:clone}{لجينيت} واحد.
\textbf{25.} يمتلئ الحقل بعد 2.3 سنة؛ ويعطي مرستيم واحد
$1.6\times 10^{6}$ نبات في تسعة أشهر؛ وعند $c = 0.8$ تستقر المتوالدة
عذريًا عند $p^* = 0.625$، وتبقي الجنسيات \qty{37.5}{\%}.
\end{solution}
